\subsection{等特征局部环}

设 A 是一个环。回忆（A，V，p. 2），当 A 含有子域时，A 的特征有定义。当且仅当 A 含有一个同构于 Q 的子域时，该特征等于 0；当且仅当有 \(p1_{A} = 0\) 时，该特征等于素数 p。若 A 的特征有定义，且 \(f: A \to B\) 是非零环同态，则 B 的特征也有定义，并且等于 A 的特征。

设 A 是一个局部环，其极大理想为 m，剩余域为 k。

\begin{enumerate}
\def\labelenumi{\alph{enumi})}
\item
  假设 \(k\) 的特征为 0。则 A 含有一个域，且 A 的特征等于 0。事实上，从 Z 到 A 的典范同态是单射，并且对于每个非零整数 \(n\)，\(n1_{\mathbf{A}}\) 在 A 中可逆，因为它不属于 m。
\item
  假设 \(k\) 的特征为 \(p \neq 0\)。则 A 含有一个域，当且仅当 \(p1_{A} = 0\)。在此情形下，A 的特征等于 \(p\)。
\end{enumerate}

设 A 是一个整局部环，其分式域为 K，剩余域为 k。

\(a^{\prime})\) 环 A 含有子域，当且仅当 \(k\) 与 K 的特征相等。在此情形下，A 的特征等于 \(k\) 和 K 的特征，并称 A 为等特征局部环。

\(b^{\prime})\) 假设域 \(k\) 和 \(\mathbf{K}\) 的特征不同。则存在素数 \(p\)，使 \(k\) 的特征为 \(p\)。由于有 \(q1_{\mathbf{A}} \neq 0\)，对每个满足 \(q \neq p\) 的素数 q 均如此，域 \(\mathbf{K}\) 的特征为 0。于是称 A 为非等特征局部环。

